% from ega4-10-fr.tex
\oldpage[IV-3]{95}
\section{Préschémas de Jacobson}
\label{section:IV.10-fr}

On a déjà eu l'occasion d'observer (5.2.5) que même les préschémas
\emph{excellents} (7.8.5) ne se comportent pas toujours comme les
«~variétés~» de la Géométrie algébrique classique, notamment en ce qui
concerne les questions de dimension~; c'est ainsi que si $X$ est le spectre
d'un anneau de valuation discrète complet, l'ensemble des points fermés de
$X$ (réduit à un seul élément) n'est pas partout dense dans $X$, et que son
complémentaire (réduit à un seul élément) est un ouvert partout dense dans
$X$ mais dont la dimension est nulle, donc $<\dim(X)$. On examine dans ce
paragraphe des conditions générales moyennant lesquelles de tels phénomènes
ne se présentent pas~; il en résulte un comportement plus satisfaisant à
certains égards pour les relations entre dimension, codimension, profondeur
et coprofondeur dans ces préschémas (10.6 et 10.8). En outre, dans les
préschémas considérés, le fait que l'ensemble des points fermés soit partout
dense (et même «~très dense~» (10.1.3)) permet de ne considérer que ces points
dans beaucoup de démonstrations~; on rejoint ainsi le point de vue classique
des «~variétés algébriques~» qui, de notre point de vue, sont les ensembles
de points fermés des \emph{préschémas algébriques sur un corps}, et on fait le
lien entre le langage des schémas et celui des «~variétés~» ou «~espaces
algébriques~» de Serre (10.9 et 10.10).

\subsection{Parties très denses d'un espace topologique}
\label{subsection:IV.10.1-fr}

\begin{env}[10.1.1]
\label{IV.10.1.1-fr}
On dit qu'une partie d'un espace topologique $X$ est
\emph{quasi-constructible} si elle est réunion finie de parties localement
fermées de $X$. On dit qu'une partie $T$ de $X$ est \emph{localement
quasi-constructible} si pour tout $x\in X$, il existe un voisinage ouvert $V$
de $x$ tel que $T\cap V$ soit quasi-constructible dans $V$. Il est clair que
toute partie quasi-constructible de $X$ est localement quasi-constructible~:
la réciproque est vraie si $X$ est quasi-compact. Le raisonnement de
($\mathbf 0_{\mathrm{III}}$, 9.1.3), où l'on supprime le mot
«~rétrocompact~», montre que l'ensemble des parties quasi-constructibles
(resp. localement quasi-constructibles) de $X$, que nous désignerons par
$\mathfrak{Qc}(X)$ (resp. $\mathfrak{Lqc}(X)$) est stable par réunion finie,
intersection finie et passage aux complémentaires. Si $f:X\to Y$ est une
application continue, il résulte aussitôt de ces définitions que, pour toute
partie quasi-constructible (resp. localement quasi-constructible) $Z$ de $Y$,
$f^{-1}(Z)$ est quasi-constructible (resp. localement quasi-constructible)
dans $X$.

Les parties constructibles (resp. localement constructibles) de $X$ sont
évidemment quasi-constructibles (resp. localement quasi-constructibles)~; la
réciproque est vraie lorsque $X$ est noethérien (resp. localement
noethérien).

Dans ce qui suit, nous désignerons respectivement par $\mathfrak O(X)$,
$\mathfrak F(X)$, $\mathfrak{Lf}(X)$, $\mathfrak C(X)$, $\mathfrak{Lc}(X)$
l'ensemble des parties de $X$ qui sont respectivement ouvertes, fermées,
localement fermées, constructibles, localement constructibles.
\end{env}

